\documentclass{article}
\usepackage{fleqn}
\usepackage{epsf}
\usepackage{aima2e-slides}

\begin{document}

\begin{huge}
\titleslide{Robot học (Robotics)}{Chương 25}

\sf

%%%%%%%%%%%% Slide %%%%%%%%%%%%%%%%%%%%%%%%%%%%%%%%%%%%%%%%%%%%%%%%%%%
\heading{Phác thảo}

Robot, Cơ quan tác động và Cảm biến

Bản địa hóa và lập bản đồ

Lập kế hoạch chuyển động

Điều khiển động cơ

%%%%%%%%%%%% Slide %%%%%%%%%%%%%%%%%%%%%%%%%%%%%%%%%%%%%%%%%%%%%%%%%%%
\heading{Robot di động}

\centerline{ \epsffile{figures/raibert-1leg.eps} \qquad \epsffile{figures/robocup.eps} }

\centerline{ \epsffile{figures/raibert-1leg.eps} \qquad \epsffile{figures/robocup.eps} }


%%%%%%%%%%%% Slide %%%%%%%%%%%%%%%%%%%%%%%%%%%%%%%%%%%%%%%%%%%%%%%%%%%
\heading{Bộ điều khiển }

\vspace*{0.2in}

\epsfxsize=0.7\textwidth
\fig{figures/stanford-arm.eps}

Cấu hình của robot được xác định bởi 6 số\al
$\implies$ 6 \defn{bậc tự do} (DOF)

6 là số lượng tối thiểu cần thiết để định vị bộ phận tác động cuối một cách tùy ý.\\
Đối với các hệ thống động lực, hãy thêm vận tốc cho mỗi DOF.


%%%%%%%%%%%% Slide %%%%%%%%%%%%%%%%%%%%%%%%%%%%%%%%%%%%%%%%%%%%%%%%%%%
\heading{Robot không ba chiều}

\vspace*{0.2in}

\epsfxsize=0.4\textwidth
\fig{figures/car-like.eps}

Một chiếc ô tô có nhiều DOF (3) hơn điều khiển (2), nên \defn{non-holonomic};\\
nói chung không thể chuyển đổi giữa hai cấu hình cực kỳ gần nhau


%%%%%%%%%%%% Slide %%%%%%%%%%%%%%%%%%%%%%%%%%%%%%%%%%%%%%%%%%%%%%%%%%%
\heading{Cảm biến}

\defn{Công cụ tìm phạm vi}: sonar (trên cạn, dưới nước), công cụ tìm phạm vi laser, radar (máy bay),\epsffile{figures/robot0124bw.eps}
cảm biến xúc giác, GPS

\vspace*{0.2in}

\centerline{ \epsffile{figures/sick.eps} \qquad \epsffile{figures/robot0124bw.eps} }

\defn{Cảm biến hình ảnh}: camera (hình ảnh, hồng ngoại)\\
\defn{ Cảm biến bản quyền }: bộ giải mã trục (khớp, bánh xe), cảm biến quán tính, \\
cảm biến lực, cảm biến mô-men xoắn


%%%%%%%%%%%% Slide %%%%%%%%%%%%%%%%%%%%%%%%%%%%%%%%%%%%%%%%%%%%%%%%%%%
\heading{Bản địa hóa---Tôi đang ở đâu?}

Tính toán vị trí và hướng hiện tại (\defn{pose}) dựa trên các quan sát:

\vspace*{0.2in}

\epsfxsize=0.7\textwidth
\fig{figures/robotics-ddn.eps}

%%%%%%%%%%%% Slide %%%%%%%%%%%%%%%%%%%%%%%%%%%%%%%%%%%%%%%%%%%%%%%%%%%
\heading{Tiếp theo bản địa hóa}

\vspace*{0.2in}

\twofig{figures/robotics-pic2.eps}{figures/range-scan-model.eps}

Giả sử nhiễu Gaussian trong dự đoán chuyển động, đo phạm vi cảm biến

%%%%%%%%%%%% Slide %%%%%%%%%%%%%%%%%%%%%%%%%%%%%%%%%%%%%%%%%%%%%%%%%%%
\heading{Tiếp theo bản địa hóa}

Có thể sử dụng tính năng lọc hạt để tạo ra ước tính vị trí gần đúng

\vspace*{0.2in}

\threefig{figures/first.ps}{figures/second.ps}{figures/third.ps}

%%%%%%%%%%%% Slide %%%%%%%%%%%%%%%%%%%%%%%%%%%%%%%%%%%%%%%%%%%%%%%%%%%
\heading{Tiếp theo bản địa hóa}

Cũng có thể sử dụng bộ lọc Kalman mở rộng \defn{} cho các trường hợp đơn giản:

\vspace*{0.2in}

\epsfxsize=0.75\textwidth
\fig{figures/robotics-pic6.eps}

Giả sử rằng các điểm mốc là {\em có thể nhận dạng được}---nếu không,
sau là đa phương thức





%%%%%%%%%%%% Slide %%%%%%%%%%%%%%%%%%%%%%%%%%%%%%%%%%%%%%%%%%%%%%%%%%%
\heading{Ánh xạ}

Bản địa hóa: bản đồ đã cho và các mốc quan sát, cập nhật phân bổ tư thế

Lập bản đồ: tư thế đã cho và các mốc quan sát, cập nhật phân bổ bản đồ

SLAM: đưa ra các mốc quan sát, cập nhật tư thế và phân bổ bản đồ

Công thức xác suất của SLAM: \al
 thêm các vị trí mốc $L_1,\ldots,L_k$ vào vectơ trạng thái,\al
 tiến hành như đối với việc bản địa hóa

%%%%%%%%%%%% Slide %%%%%%%%%%%%%%%%%%%%%%%%%%%%%%%%%%%%%%%%%%%%%%%%%%%
\heading{Tiếp theo ánh xạ }

\qquad

\centerline{\epsfysize=0.37\textheight\epsffile{figures/arena033.ps}\qquad\epsfysize=0.37\textheight\epsffile{figures/arena034.ps}}
\centerline{\epsfysize=0.37\textheight\epsffile{figures/arena033.ps}\qquad\epsfysize=0.37\textheight\epsffile{figures/arena034.ps}}

%%%%%%%%%%%% Slide %%%%%%%%%%%%%%%%%%%%%%%%%%%%%%%%%%%%%%%%%%%%%%%%%%%
\heading{Ví dụ về bản đồ 3D}

\vspace*{0.2in}

\twofig{figures/mine-robot.eps}{figures/mine-data.eps}

%%%%%%%%%%%% Slide %%%%%%%%%%%%%%%%%%%%%%%%%%%%%%%%%%%%%%%%%%%%%%%%%%%
\heading{Lập kế hoạch chuyển động}

Ý tưởng: sơ đồ trong không gian cấu hình \defn{} được xác định bởi DOF của robot

\epsffile{figures/armExampleConfSpace.eps}

\centerline{ \epsffile{figures/armExampleWorkSpace.eps} \qquad \epsffile{figures/armExampleConfSpace.eps} }

Giải là quỹ đạo điểm trong không gian C tự do

%%%%%%%%%%%% Slide %%%%%%%%%%%%%%%%%%%%%%%%%%%%%%%%%%%%%%%%%%%%%%%%%%%
\heading{Quy hoạch không gian cấu hình}

Vấn đề cơ bản: $\infty^d$ khẳng định! Chuyển đổi sang không gian trạng thái \emph{hữu hạn}.

\defn{Phân hủy tế bào}:\al
  chia không gian thành các ô \defn{đơn giản}, \al
  mỗi trong số đó có thể được duyệt qua một cách “dễ dàng” (ví dụ: lồi)

\defn{ Khung xương }: \al
  xác định số lượng hữu hạn các điểm/đường được kết nối dễ dàng\al
  tạo thành một biểu đồ sao cho hai điểm bất kỳ được kết nối\al
  bởi một đường dẫn trên đồ thị


%%%%%%%%%%%% Slide %%%%%%%%%%%%%%%%%%%%%%%%%%%%%%%%%%%%%%%%%%%%%%%%%%%
\heading{Ví dụ về phân hủy ô}

\qquad

\centerline{ \epsffile{figures/armDPwithoutPotentialWorkspaceCoarse.eps} \qquad \epsffile{figures/armDPwithoutPotentialCoarse.eps} }

Sự cố: có thể không có đường dẫn trong các ô không gian trống thuần túy\\
Giải pháp: phân rã đệ quy các ô hỗn hợp (tự do + chướng ngại vật)


%%%%%%%%%%%% Slide %%%%%%%%%%%%%%%%%%%%%%%%%%%%%%%%%%%%%%%%%%%%%%%%%%%
\heading{Skeletonization: Sơ đồ Voronoi}

Sơ đồ Voronoi: quỹ tích các điểm cách đều chướng ngại vật

\vspace*{0.2in}

\epsfxsize=0.44\textwidth
\fig{figures/armVoronoi.eps}

Vấn đề: không mở rộng tốt sang các kích thước cao hơn

%%%%%%%%%%%% Slide %%%%%%%%%%%%%%%%%%%%%%%%%%%%%%%%%%%%%%%%%%%%%%%%%%%
\heading{Skeletonization: Lộ trình xác suất}

Lộ trình xác suất được tạo bằng cách tạo các điểm ngẫu nhiên trong
C-space và giữ chúng trong không gian tự do; tạo biểu đồ bằng cách tham gia
cặp theo đường thẳng

\vspace*{0.2in}

\epsfxsize=0.44\textwidth
\fig{figures/armRoadmap.eps}

Vấn đề: cần tạo đủ điểm để đảm bảo rằng mọi
cặp điểm bắt đầu/mục tiêu được kết nối thông qua biểu đồ

%%%%%%%%%%%% Slide %%%%%%%%%%%%%%%%%%%%%%%%%%%%%%%%%%%%%%%%%%%%%%%%%%%
\heading{Điều khiển động cơ}

Có thể xem vấn đề điều khiển động cơ dưới dạng vấn đề tìm kiếm\\
trong không gian trạng thái \note{dynamic} thay vì \note{kinematic}:\al
-- không gian trạng thái được xác định bởi $x_1,x_2,\ldots,\dot{x_1},\dot{x_2},\ldots$\al
-- liên tục, đa chiều (Sarcos hình người: 162 chiều)

Kiểm soát xác định: nhiều vấn đề có thể giải quyết chính xác\\
đặc biệt~nếu tuyến tính, ít chiều, được biết chính xác, có thể quan sát được

Luật điều khiển \defn{đơn giản} có hiệu lực đối với các chuyển động cụ thể

Stochastic \defn{điều khiển tối ưu}: rất ít vấn đề có thể giải được chính xác\al
$\implies$ phương pháp gần đúng/thích ứng




%%%%%%%%%%%% Slide %%%%%%%%%%%%%%%%%%%%%%%%%%%%%%%%%%%%%%%%%%%%%%%%%%%
\heading{Điều khiển động cơ sinh học}

Hệ thống điều khiển động cơ được đặc trưng bởi sự dư thừa lớn

Vô số quỹ đạo đạt được bất kỳ nhiệm vụ nhất định nào

Ví dụ: cánh tay 3 liên kết di chuyển trong mặt phẳng ném vào mục tiêu\nl
  bộ điều khiển 12 tham số đơn giản, một bậc tự do tại mục tiêu\nl
  Không gian liên tục 11 chiều của bộ điều khiển tối ưu

Ý tưởng: nếu cánh tay ồn ào, chỉ có chính sách tối ưu `` một ''
giảm thiểu sai sót ở mục tiêu

Tức là, khả năng chịu tiếng ồn có thể giải thích hành vi vận động thực tế

Harris \& Wolpert ({\em Nature}, 1998): nhiễu phụ thuộc vào tín hiệu\\
giải thích hồ sơ vận tốc mắt chớp một cách hoàn hảo

%%%%%%%%%%%% Slide %%%%%%%%%%%%%%%%%%%%%%%%%%%%%%%%%%%%%%%%%%%%%%%%%%%
\heading{Cài đặt}

Giả sử một bộ điều khiển có các tham số điều khiển ``dự kiến'' $\theta_0$\\
bị hỏng do nhiễu, tạo ra $\theta$ được rút ra từ $P_{\theta_0}$

Đầu ra (ví dụ: khoảng cách từ mục tiêu) $y = F(\theta)$; 


\epsfxsize=0.7\maxfigwidth
\fig{figures/arm-setup.ps}


%%%%%%%%%%%% Slide %%%%%%%%%%%%%%%%%%%%%%%%%%%%%%%%%%%%%%%%%%%%%%%%%%%
\heading{Thuật toán học đơn giản: Độ dốc ngẫu nhiên}

Giảm thiểu $E_{\theta}[y^2]$ bằng cách giảm độ dốc:
\begin{eqnarray*}
\nabla_{\theta_0} E_{\theta}[y^2]
  &=&  \nabla_{\theta_0} \int P_{\theta_0}(\theta)F(\theta)^2 d\theta \\
  &=&  \int \frac{\nabla_{\theta_0} P_{\theta_0}(\theta)}{P_{\theta_0}(\theta)}
            F(\theta)^2 P_{\theta_0}(\theta)d\theta \\
  &=&  E_{\theta}[\frac{\nabla_{\theta_0} P_{\theta_0}(\theta)}
                       {P_{\theta_0}(\theta)}
                  y^2]
\end{eqnarray*}
Cho mẫu $(\theta_j,y_j)$, $j=1,\ldots,N$, ta có
\[
  \hat{\nabla_{\theta_0} E_{\theta}[y^2]} =
  \frac{1}{N}\sum_{j=1}^N
     \frac{\nabla_{\theta_0} P_{\theta_0}(\theta_j)}
                       {P_{\theta_0}(\theta_j)}
                  y_j^2
\]
Đối với nhiễu Gaussian có hiệp phương sai $\Sigma$, tức là
$P_{\theta_0}(\theta) = N(\theta_0,\Sigma)$, chúng tôi thu được
\[
  \hat{\nabla_{\theta_0} E_{\theta}[y^2]} =
  \frac{1}{N}\sum_{j=1}^N \Sigma^{-1}(\theta_j - \theta_0) y_j^2
\]


%%%%%%%%%%%% Slide %%%%%%%%%%%%%%%%%%%%%%%%%%%%%%%%%%%%%%%%%%%%%%%%%%%
\heading{Thuật toán đang làm gì}

\vspace*{0.3in}

\epsfxsize=0.7\maxfigwidth
\fig{figures/control-noise-distort.ps}


%%%%%%%%%%%% Slide %%%%%%%%%%%%%%%%%%%%%%%%%%%%%%%%%%%%%%%%%%%%%%%%%%%
\heading{Kết quả cho bộ điều khiển 2--D}

\vspace*{0.3in}

\epsfxsize=0.75\textwidth
\centerline{ \epsffile{figures/progress3-theta.ps} }

%%%%%%%%%%%% Slide %%%%%%%%%%%%%%%%%%%%%%%%%%%%%%%%%%%%%%%%%%%%%%%%%%%
\heading{Kết quả cho bộ điều khiển 2--D}

\vspace*{0.3in}

\epsfxsize=0.75\textwidth
\centerline{ \epsffile{figures/progress3-theta-zoom.ps} }



%%%%%%%%%%%% Slide %%%%%%%%%%%%%%%%%%%%%%%%%%%%%%%%%%%%%%%%%%%%%%%%%%%
\heading{Kết quả cho bộ điều khiển 2--D}

\vspace*{0.3in}

\epsfxsize=0.75\textwidth
\centerline{ \epsffile{figures/progress3-score.ps} }


%%%%%%%%%%%% Slide %%%%%%%%%%%%%%%%%%%%%%%%%%%%%%%%%%%%%%%%%%%%%%%%%%%
\heading{Tóm tắt}

Cao su rơi xuống đường

Robot di động và người thao tác

Bậc tự do xác định cấu hình robot

Bản địa hóa và ánh xạ dưới dạng các vấn đề suy luận xác suất\al
  (yêu cầu mô hình cảm biến và chuyển động tốt)

Lập kế hoạch chuyển động trong không gian cấu hình\al
  yêu cầu một số phương pháp để hoàn thiện



\end{huge}
\end{document}


